% =====================================================================
%  CHƯƠNG 4: KẾT QUẢ THỰC NGHIỆM
% =====================================================================
\chapter{Kết quả thực nghiệm}
\label{chap:exp}

\section{Môi trường và cấu hình}

Toàn bộ thực nghiệm được triển khai trên môi trường điện toán đám mây
\textbf{Google Colab} với bộ xử lý đồ họa (GPU) gia tốc. Để tối ưu vùng nhớ và tốc
độ, đồ án sử dụng kỹ thuật \textbf{tính toán độ chính xác hỗn hợp} (mixed precision)
thông qua \texttt{torch.amp} kết hợp \texttt{GradScaler}. Bảng~\ref{tab:hyper} tổng
hợp các siêu tham số chính.

\begin{table}[H]
\centering
\caption{Cấu hình môi trường và siêu tham số huấn luyện}
\label{tab:hyper}
\renewcommand{\arraystretch}{1.3}
\begin{tabularx}{\textwidth}{@{}l X@{}}
\toprule
\textbf{Thành phần} & \textbf{Cấu hình} \\
\midrule
Bộ mã hóa văn bản (tháp Item) & \texttt{all-MiniLM-L6-v2} ($6$ tầng, $384$ chiều), đóng băng \\
Tháp người dùng & MLP $384{\to}512{\to}384$, BatchNorm, ReLU, Dropout $0{,}5$ \\
Hàm mất mát & InfoNCE (in-batch negatives), temperature $\tau = 0{,}1$ \\
Bộ tối ưu & AdamW, learning rate $10^{-4}$, weight decay $0{,}05$ \\
Batch size & $256$ \\
Số epoch & $5$ (early stopping theo val loss) \\
Độ chính xác & Mixed precision (FP16) trên GPU \\
Nạp dữ liệu & Parquet + DataLoader $2$ worker, bất đồng bộ \\
\bottomrule
\end{tabularx}
\end{table}

Quy mô dữ liệu: khoảng $984.000$ tương tác huấn luyện ($3.844$ batch mỗi epoch),
$146.396$ người dùng và hơn $17.000$ sản phẩm. Bảng~\ref{tab:time} phân rã thời gian
một epoch: phần lớn dành cho huấn luyện ($\sim 38$ phút), phần còn lại cho tính mất
mát kiểm định và đánh giá toàn kho ($\sim 11$ phút).

\begin{table}[H]
\centering
\caption{Phân rã thời gian một epoch (tốc độ $\sim 1{,}7$ batch/giây)}
\label{tab:time}
\renewcommand{\arraystretch}{1.25}
\begin{tabularx}{0.8\textwidth}{@{}X r@{}}
\toprule
\textbf{Giai đoạn} & \textbf{Thời gian} \\
\midrule
Huấn luyện ($3.844$ batch)            & $\approx 38$ phút \\
Tính mất mát kiểm định ($572$ batch)  & $\approx 5{,}7$ phút \\
Mã hóa kho item + tính Recall/NDCG    & $\approx 5{,}4$ phút \\
\midrule
\textbf{Tổng mỗi epoch}               & $\approx 49$ phút \\
\bottomrule
\end{tabularx}
\end{table}

\section{Chi tiết vòng lặp huấn luyện}

Vòng lặp huấn luyện kết hợp nhiều kỹ thuật để tối ưu tốc độ và bộ nhớ. Mỗi batch
trải qua các bước: (1) mã hóa sản phẩm mục tiêu bằng tháp Item (không gradient);
(2) đưa vector lịch sử qua tháp User; (3) tính InfoNCE; (4) lan truyền ngược. Toàn
bộ phép tính tiến (forward) được bọc trong \texttt{autocast} để dùng FP16, còn
\texttt{GradScaler} nhân tỷ lệ mất mát nhằm tránh tràn số (underflow) khi gradient
quá nhỏ ở độ chính xác thấp.

\begin{lstlisting}[language=Python, caption=Vòng lặp huấn luyện với mixed precision và early stopping]
scaler = torch.amp.GradScaler('cuda')
best_val_loss = float('inf'); patience_counter = 0

for epoch in range(EPOCHS):
    model.train(); total_loss = 0
    for batch in train_loader:
        item_tokens = {k: v.to(device) for k, v in batch["item_tokens"].items()}
        optimizer.zero_grad()
        with torch.amp.autocast('cuda'):          # mixed precision FP16
            item_out = model.forward_item(item_tokens)        # thap dong bang
            user_out = model.forward_user(batch["user_vectors"].to(device))
            loss = loss_fn(user_out, item_out)                # InfoNCE
        scaler.scale(loss).backward()              # scale tranh underflow
        scaler.step(optimizer); scaler.update()
        total_loss += loss.item()

    val_loss = compute_loss(model, val_loader, loss_fn, device)
    if val_loss < best_val_loss:                   # early stopping theo val loss
        best_val_loss = val_loss; patience_counter = 0
        torch.save(model.state_dict(), best_model_path)
    else:
        patience_counter += 1
        if patience_counter >= PATIENCE: break
\end{lstlisting}

Cơ chế \textbf{early stopping} lưu lại trọng số tốt nhất mỗi khi mất mát kiểm định
giảm và dừng huấn luyện nếu không cải thiện sau \texttt{PATIENCE} epoch -- bảo đảm
mô hình cuối cùng là phiên bản tổng quát hóa tốt nhất, tránh lãng phí tính toán.

\section{Thiết kế nạp dữ liệu theo batch}

Việc nạp dữ liệu được tổ chức qua lớp \texttt{Dataset} trả về vector người dùng (đã
tính sẵn) cùng văn bản sản phẩm mục tiêu, và một hàm \texttt{collate} gom batch:
chỉ \textit{văn bản sản phẩm} mới cần tokenize tại thời điểm tạo batch, còn vector
người dùng đã được mã hóa trước nên chỉ việc xếp chồng -- giảm đáng kể khối lượng
tính toán lặp lại. \texttt{DataLoader} dùng $2$ worker và \texttt{pin\_memory} để
nạp bất đồng bộ, che giấu độ trễ I/O sau thời gian tính toán của GPU.

\begin{lstlisting}[language=Python, caption=Hàm collate chỉ tokenize nhánh sản phẩm]
def collate_vector_fn(batch):
    user_vectors = torch.stack([b["user_vector"] for b in batch])  # [B, 384]
    item_texts = [b["item_text"] for b in batch]
    item_tokens = tokenizer(item_texts, padding=True, truncation=True,
                            max_length=128, return_tensors="pt")
    return {"user_vectors": user_vectors, "item_tokens": item_tokens,
            "target": [b["target"] for b in batch]}
\end{lstlisting}

\section{Giao thức đánh giá toàn kho (Full-catalog)}

Một đóng góp quan trọng về mặt phương pháp là cách đánh giá. Nhiều nghiên cứu chỉ
đánh giá trên một nhóm nhỏ mẫu âm lấy ngẫu nhiên (in-batch / sampled metrics), dẫn
tới chỉ số bị \textit{thổi phồng} và không phản ánh độ khó thực tế. Đồ án triệt tiêu
hoàn toàn thiên lệch này bằng giao thức \textbf{full-catalog}:

\begin{enumerate}[leftmargin=2em, itemsep=2pt]
    \item Mã hóa trước \textbf{toàn bộ} hơn $17.000$ sản phẩm thành ma trận vector
    sản phẩm.
    \item Với mỗi người dùng, tính độ tương đồng với \textit{tất cả} sản phẩm trong
    kho, lấy Top-$50$.
    \item Đo Recall@K và NDCG@K bằng cách kiểm tra sản phẩm mục tiêu thật có nằm
    trong Top-$K$ hay không.
\end{enumerate}

Việc tính toán được chia theo cụm người dùng ($1.000$ người mỗi lần) để tránh tràn
bộ nhớ GPU. Đây là cách đánh giá \textit{khắt khe và trung thực} nhất, vì người dùng
phải cạnh tranh với toàn bộ kho sản phẩm.

Về mặt kỹ thuật, nếu nhân trực tiếp ma trận người dùng ($146.396 \times 384$) với
ma trận sản phẩm ($17.000 \times 384$) sẽ tạo ra ma trận điểm số khổng lồ vượt quá
VRAM. Giải pháp là \textbf{xử lý theo cụm}: kho sản phẩm được mã hóa trước thành một
ma trận cố định; sau đó với mỗi cụm $1.000$ người dùng, ta tính điểm số, lấy
\texttt{topk} rồi giải phóng bộ nhớ trước khi sang cụm tiếp theo.

\begin{lstlisting}[language=Python, caption=Đánh giá toàn kho theo cụm để tránh tràn VRAM]
# Ma hoa truoc TOAN BO kho item mot lan
all_item_emb = encode_all_items(item_ids, model)        # [N_items, 384]

for i in range(0, num_users, 1000):                      # xu ly 1000 user/lan
    user_chunk = all_user_emb[i:i+1000].to(device)
    scores = torch.matmul(user_chunk, all_item_emb.T)    # [1000, N_items]
    _, topk = torch.topk(scores, k=50, dim=1)            # lay Top-50
    hits = (topk == target_chunk.view(-1, 1))            # vi tri trung
    recall10 += hits[:, :10].any(dim=1).float().sum().item()
    # ... cong don NDCG theo thu hang ...
    del scores, topk; torch.cuda.empty_cache()           # giai phong VRAM
\end{lstlisting}

Nhờ chiến lược này, hệ thống có thể đánh giá trên toàn bộ $146.396$ người dùng $\times$
$17.000$ sản phẩm mà vẫn nằm trong giới hạn bộ nhớ của một GPU phổ thông.

\section{Kết quả huấn luyện}

Hình~\ref{fig:loss} thể hiện đường cong mất mát qua $5$ epoch. Mất mát huấn luyện
(InfoNCE) giảm rõ từ $3{,}62$ xuống $3{,}31$, trong khi mất mát kiểm định
\textit{giảm đều đặn} từ $5{,}19$ xuống $5{,}16$ -- cho thấy mô hình hội tụ ổn định
và \textbf{không bị quá khớp}. Cơ chế early stopping lưu lại mô hình tốt nhất tại
mỗi lần val loss giảm; mô hình đạt trạng thái ổn định tại epoch thứ $5$.

\begin{figure}[H]
\centering
\includegraphics[width=0.78\textwidth]{loss_curve.png}
\caption{Đường cong mất mát InfoNCE trên tập huấn luyện và kiểm định}
\label{fig:loss}
\end{figure}

\section{Kết quả truy hồi trên toàn kho}

Bảng~\ref{tab:results} trình bày kết quả Recall@K và NDCG@K theo từng epoch trên tập
kiểm định, trực quan hóa ở Hình~\ref{fig:metrics}. Tại epoch $5$, mô hình đạt
Recall@10 $= 0{,}0142$, Recall@50 $= 0{,}0436$, NDCG@10 $= 0{,}0062$ và NDCG@50
$= 0{,}0125$.

\begin{table}[H]
\centering
\caption{Kết quả truy hồi toàn kho theo từng epoch (tập kiểm định)}
\label{tab:results}
\renewcommand{\arraystretch}{1.3}
\begin{tabularx}{\textwidth}{@{}c X X X X X@{}}
\toprule
\textbf{Epoch} & \textbf{Train loss} & \textbf{Val loss} & \textbf{Recall@10} &
\textbf{Recall@50} & \textbf{NDCG@10} \\
\midrule
1 & $3{,}6168$ & $5{,}1875$ & $0{,}0140$ & $0{,}0431$ & $0{,}0061$ \\
2 & $3{,}3495$ & $5{,}1790$ & $0{,}0140$ & $0{,}0435$ & $0{,}0061$ \\
3 & $3{,}3229$ & $5{,}1755$ & $0{,}0141$ & $0{,}0429$ & $0{,}0061$ \\
4 & $3{,}3138$ & $5{,}1692$ & $0{,}0141$ & $0{,}0434$ & $0{,}0061$ \\
\textbf{5} & $\mathbf{3{,}3076}$ & $\mathbf{5{,}1649}$ & $\mathbf{0{,}0142}$ &
$\mathbf{0{,}0436}$ & $\mathbf{0{,}0062}$ \\
\bottomrule
\end{tabularx}
\end{table}

\begin{figure}[H]
\centering
\includegraphics[width=\textwidth]{metrics_curve.png}
\caption{Recall@K (trái) và NDCG@K (phải) trên tập kiểm định qua các epoch}
\label{fig:metrics}
\end{figure}

\subsection{So sánh với baseline ngẫu nhiên}

Các con số tuyệt đối nói trên thoạt nhìn có vẻ nhỏ, nhưng cần đặt trong bối cảnh
\textbf{kho hơn $17.000$ sản phẩm}. Nếu chọn ngẫu nhiên, xác suất trúng sản phẩm
mục tiêu trong Top-$10$ chỉ là $10/17000 \approx 0{,}00059$. Như vậy, mô hình đề
xuất đạt kết quả \textbf{vượt trội hơn $24$ lần} ở Top-$10$ và hơn $14$ lần ở
Top-$50$ so với chọn ngẫu nhiên (Hình~\ref{fig:uplift}, lưu ý trục tung thang log).

\begin{table}[H]
\centering
\caption{Bội số cải thiện của mô hình so với chọn ngẫu nhiên}
\label{tab:uplift}
\renewcommand{\arraystretch}{1.3}
\begin{tabularx}{0.9\textwidth}{@{}X r r r@{}}
\toprule
\textbf{Chỉ số} & \textbf{Ngẫu nhiên} & \textbf{Two-Tower} & \textbf{Bội số} \\
\midrule
Recall@10 & $0{,}00059$ & $0{,}0142$ & $\approx 24{,}1\times$ \\
Recall@50 & $0{,}00294$ & $0{,}0436$ & $\approx 14{,}8\times$ \\
\bottomrule
\end{tabularx}
\end{table}

\begin{figure}[H]
\centering
\includegraphics[width=0.72\textwidth]{uplift_vs_random.png}
\caption{Hiệu năng truy hồi của mô hình so với baseline chọn ngẫu nhiên}
\label{fig:uplift}
\end{figure}

\section{Ví dụ gợi ý định tính}

Để minh họa trực quan năng lực của mô hình, Bảng~\ref{tab:qual} trình bày Top-$5$
sản phẩm được gợi ý cho một người dùng tiêu biểu có lịch sử thiên về phần mềm bảo
mật và tiện ích hệ thống. Có thể thấy mô hình nắm bắt đúng "khẩu vị" ngữ nghĩa: các
gợi ý đều thuộc nhóm bảo mật/tiện ích, dù người dùng chưa từng tương tác trực tiếp
với chúng -- minh chứng cho khả năng khái quát hóa qua không gian vector ngữ nghĩa.

\begin{table}[H]
\centering
\caption{Ví dụ Top-$5$ gợi ý dựa trên độ tương đồng vector}
\label{tab:qual}
\renewcommand{\arraystretch}{1.3}
\begin{tabularx}{\textwidth}{@{}c X c@{}}
\toprule
\textbf{Hạng} & \textbf{Sản phẩm gợi ý} & \textbf{Cosine} \\
\midrule
1 & Norton 360 Deluxe -- Antivirus \& VPN            & $0{,}71$ \\
2 & Malwarebytes Premium -- Anti-Malware             & $0{,}68$ \\
3 & McAfee Total Protection                          & $0{,}66$ \\
4 & CCleaner Professional -- System Utility          & $0{,}63$ \\
5 & Acronis Cyber Protect -- Backup \& Security       & $0{,}61$ \\
\bottomrule
\end{tabularx}
\end{table}

\section{Phân tích không gian biểu diễn}

Một góc nhìn hữu ích là so sánh mất mát với \textit{mức ngẫu nhiên}. Với batch
$256$, hàm InfoNCE của một mô hình đoán ngẫu nhiên có giá trị $\ln(256) \approx
5{,}55$. Như vậy, mất mát kiểm định $5{,}16$ chỉ \textit{nhỉnh hơn ngẫu nhiên một
chút}, trong khi mất mát huấn luyện $3{,}31$ thấp hơn rõ rệt. Khoảng cách này cho
biết: mô hình đã sắp xếp tốt không gian cho các tương tác trong tập huấn luyện, nhưng
\textbf{khái quát hóa sang người dùng kiểm định còn hạn chế} -- điều dễ hiểu khi mục
tiêu kiểm định (sản phẩm kề cuối theo thời gian) là một bài toán khó và tháp sản phẩm
bị đóng băng nên không gian biểu diễn không được "uốn nắn" theo dữ liệu tương tác.

Điều quan trọng là mất mát kiểm định vẫn \textit{giảm đơn điệu} qua cả $5$ epoch:
mô hình không rơi vào quá khớp theo nghĩa cổ điển (val loss tăng trở lại), mà chạm
tới một \textbf{trần khái quát hóa} do thiết kế đóng băng. Đây chính là cơ sở định
lượng cho đề xuất "mở băng" tháp sản phẩm bằng LoRA ở chương kết luận.

Về mặt định tính, khi chiếu các vector sản phẩm xuống không gian $2$ chiều (bằng
t-SNE), ta kỳ vọng các sản phẩm cùng nhóm chức năng tụ lại thành cụm
(Hình~\ref{fig:embspace}). Vector người dùng được đặt gần cụm phản ánh sở thích của
họ, và Top-$N$ gợi ý chính là các sản phẩm gần nhất xung quanh điểm đó. Cấu trúc cụm
này -- kế thừa từ tri thức ngôn ngữ của MiniLM -- là lý do mô hình hoạt động tốt cho
cả sản phẩm mới chưa có lịch sử tương tác.

\begin{figure}[H]
\centering
\includegraphics[width=0.78\textwidth]{embedding_space.png}
\caption{Minh họa khái niệm không gian biểu diễn $2$ chiều của sản phẩm và người dùng}
\label{fig:embspace}
\end{figure}

\section{Phân tích và thảo luận}

Từ các kết quả trên, có thể rút ra một số nhận định:

\begin{itemize}[leftmargin=2em, itemsep=2pt]
    \item \textbf{Mô hình học được biểu diễn có ý nghĩa:} dù chỉ huấn luyện tháp
    người dùng (tháp sản phẩm đóng băng), mô hình vẫn ánh xạ thành công lịch sử hành
    vi sang không gian sản phẩm, đạt mức cải thiện hàng chục lần so với ngẫu nhiên
    trên một kho rất lớn.
    \item \textbf{Hội tụ nhanh và ổn định:} các chỉ số gần như bão hòa ngay từ epoch
    đầu và nhích nhẹ qua các epoch sau, cho thấy bộ mã hóa MiniLM đóng băng đã cung
    cấp một không gian biểu diễn tốt; phần việc còn lại của tháp người dùng tương đối
    "dễ" và nhanh đạt giới hạn.
    \item \textbf{Trần hiệu năng đến từ tháp đóng băng:} vì vector sản phẩm cố định
    (không được tinh chỉnh theo dữ liệu tương tác), không gian biểu diễn chưa được
    "uốn nắn" theo hành vi mua sắm thực tế -- đây vừa là ưu điểm (nhẹ, nhanh, chống
    cold-start) vừa là hạn chế về độ chính xác, sẽ được bàn ở chương kết luận.
    \item \textbf{Giá trị của đánh giá full-catalog:} con số tuyệt đối khiêm tốn
    nhưng \textit{trung thực}; nếu dùng đánh giá lấy mẫu nhỏ, các chỉ số sẽ cao hơn
    nhiều nhưng không phản ánh đúng năng lực thực tế của hệ thống.
\end{itemize}
